% Escreva aqui a solução da questão 1.
\textbf{VERDADEIRA}

\textcolor{red}{ VERSÃO 1 } %Letícia

%\textbf{VERDADEIRA}

De uma das propriedades da união de conjuntos, sabemos: $A \subseteq (A \cup B)$.

Das propriedades de complementar de um conjunto, sabemos: $X \subseteq Y \implies Y^C \subseteq X^C$.

Aplicando essa propriedade à primeira inclusão acima:

$A \subseteq (A \cup B) \implies (A \cup B)^{C} \subseteq A^{C}$.
Portanto, a sentença é verdadeira.

\bigskip

\textcolor{red}{CHAVE DE PONTUAÇÃO:
\begin{itemize}
    \item  Inclusão pela união \emph{(1,5 ponto)};
    \item  Inclusão pela propriedade de complementar \emph{(1,5 ponto)}.
\end{itemize} }


\bigskip


\textcolor{red}{ VERSÃO 2} %Esther



Perceba que, por De Morgan, $(A \cup B)^C = A^C\cap B^C$ e, em particular, $(A \cup B)^C \subseteq A^C\cap B^C$. Por propriedade de interseção, $A^C\cap B^C\subseteq A^C$. Logo, por transitividade, $(A \cup B)^C \subseteq A^C.$


\bigskip


%\textcolor{red}{ VERSÃO 3} % autor




%\bigskip







\textcolor{red}{CHAVE DE PONTUAÇÃO:
\begin{itemize}
    \item  Inclusão pelo DeMorgan \emph{(1,0 ponto)};
    \item  Inclusão pela interseção \emph{(1,0 ponto)};
    \item  Inclusão pela transitividade \emph{(1,0 ponto)}.
\end{itemize} }